\writeSectionFrame{\presentationTitle}{Adapter}
\begin{frame}{Steckbrief}
	\begin{itemize}
 		\item Deutscher Name: Adapter
 		\item Englischer Name: Adapter
 		\item Gruppe: Strukturmuster
 	\end{itemize}
\end{frame}

\begin{frame}{Adapter}
	\begin{block}{Zweck}
 		Passe die Schnittstelle einer Klasse an eine andere
        von ihren Klienten erwartete Schnittstelle an. Das Adaptermuster lässt Klassen zusammenarbeiten, die wegen inkompatibler Schnittstellen sonst dazu nicht in der Lage wären.
 	\end{block}
\end{frame}

\begin{frame}{Adapterarten}
	\begin{itemize}
 		\item Klassenadapter
 		\item Objektadapter
 	\end{itemize}
\end{frame}

\begin{frame}{Allgemeines Klassendiagramm - Klassenadapter}
    \begin{tikzpicture}[scale=0.7, transform shape]
        \umlclass[x=-3,y=0]{Client}{
        }{
        }
    
        \umlclass[x=2,y=0]{Ziel}{
        }{
            +operationDesZiels()
        }

        \umlclass[x=8,y=0]{AdaptierteKlasse}{
        }{
            +operationDerAdaptiertenKlasse()
        }

        \umlclass[x=4,y=-3]{Adapter}{
        }{
            +operationDesZiels()
        }
        
        \umluniassoc{Client}{Ziel}
        
        \umlinherit{Adapter}{Ziel}
        \umlinherit{Adapter}{AdaptierteKlasse}
    \end{tikzpicture}
\end{frame}

\note{
	Der Klassenadapter erbt von der Zielklasse, also die Klasse, die ein Klient verwenden
	möchte, und der adaptierten Klasse. Das Problem wird also während der Entwurfszeit gelöst. 
	Diese Art eines Adapters ist allerdings nur mit einer Programmiersprache möglich, 
	die Mehrfachvererbung unterstützt und bringt natürlich auch die Nachteile von Mehrfachvererbung
	mit sich.  
}

\begin{frame}{Allgemeines Klassendiagramm - Objektadapter}
    \begin{tikzpicture}[scale=0.7, transform shape]
        \umlclass[x=-3,y=0]{Client}{
        }{
        }
    
        \umlclass[x=2,y=0]{Ziel}{
        }{
            +operationDesZiels()
        }

        \umlclass[x=8,y=0]{AdaptierteKlasse}{
        }{
            +operationDerAdaptiertenKlasse()
        }

        \umlclass[x=4,y=-3]{Adapter}{
        }{
            +operationDesZiels()
        }
        
        \umluniassoc{Client}{Ziel}
        
        \umlinherit{Adapter}{Ziel}
        \umluniassoc{Adapter}{AdaptierteKlasse}
    \end{tikzpicture}
\end{frame}

\begin{frame}{Klassendiagramm - Beispiel}
    \begin{tikzpicture}[scale=0.7, transform shape]
        \umlclass[x=-3,y=0]{WetterWebseite}{
        }{
        }
    
        \umlclass[x=2,y=0]{Wetterdienst}{
        }{
        }

        \umlclass[x=8,y=0]{WetterdienstXml}{
        }{
        }

        \umlclass[x=6,y=-3]{WetterdienstXmlAdapter}{
        }{
        
        }

        \umlclass[x=0,y=-3]{WetterdienstJson}{
        }{
        }

        \umlclass[x=8,y=3]{XmlText}{
        }{
        }

        \umlclass[x=0,y=-5]{JsonText}{
        }{
        }
        
        \umluniassoc{WetterWebseite}{Wetterdienst}
        \umluniassoc{WetterdienstJson}{JsonText}
        \umluniassoc{WetterdienstXml}{XmlText}
        
        \umlinherit{WetterdienstXmlAdapter}{Wetterdienst}
        \umlinherit{WetterdienstJson}{Wetterdienst}
        \umluniassoc{WetterdienstXmlAdapter}{WetterdienstXml}
    \end{tikzpicture}
\end{frame}
\begin{frame}[fragile]{Ziel}
	\begin{exampleblock}{Wetterdienst.java}
 		\begin{lstlisting}
public interface Wetterdienst {
	void empfangeDaten(JsonText wetterdaten);
	JsonText sendeDaten();
}
 		\end{lstlisting}
 	\end{exampleblock}
\end{frame}

\begin{frame}[fragile]{Konkretes Ziel}
	\begin{exampleblock}{WetterdienstJson}
        \begin{lstlisting}
public class WetterdienstJson implements Wetterdienst { 
	public void empfangeDaten(JsonText wetterdaten) {
		System.out.println("Empfange Daten");
	}

	public JsonText sendeDaten() {
		return new JsonText("{"
				+ "wetterdaten: "
				+ "{"
				+ "wetter: sonnig"
				+ "temperatur: 30"
				+ "}"
				+ "}");
	}
}
        \end{lstlisting}
 	\end{exampleblock}
\end{frame}

\begin{frame}[fragile]{AdaptierteKlasse}
	\begin{exampleblock}{WetterdienstXml.java}
 		\begin{lstlisting}
public class WetterdienstXml {
	public XmlText sendeDaten() {
		return new XmlText("<wetterdaten>"
				+ "<wetter>sonnig</wetter>"
				+ "<temperatur>30</temperatur>"
				+ "</wetterdaten>", "wetter.dtd");
	}

	public void empfangeDaten(WetterdatenDatei wetterdaten) {
		System.out.println("WetterdienstXml empfaengt Daten");
	}
} 		    
 		\end{lstlisting}
 	\end{exampleblock}
\end{frame}

\begin{frame}[fragile]{Adapter}
	\begin{exampleblock}{WetterdienstXmlAdapter.java}
 		\begin{lstlisting}
public class WetterdienstXmlAdapter implements Wetterdienst {
	private final WetterdienstXml wetterdienstXml;
	
	@Override
	public JsonText sendeDaten() {
		System.out.println("XML-Daten werden empfangen");
		XmlText wetterdatenAlsXml = wetterdienstXml.sendeDaten();
		System.out.println("XML-Daten zu JSON konvertiert");
		return konvertiereZuJson(wetterdatenAlsXml);
	}
    // ...
} 		    
 		\end{lstlisting}
 	\end{exampleblock}
\end{frame}

\begin{frame}[fragile]{Klient}
	\begin{exampleblock}{WetterWebseite.java}
 		\begin{lstlisting}
public class WetterWebseite {
	private List<Wetterdienst> wetterdienste;
	
	public WetterWebseite() {
		wetterdienste = new ArrayList<>();
		wetterdienste.add(new WetterdienstJson());
		WetterdienstXml wetterdienstXml = new WetterdienstXml();
		wetterdienste.add(new WetterdienstXmlAdapter(wetterdienstXml));
	}
	
	public void zeigeWetterdatenAn() {
		for (Wetterdienst wetterdienst: wetterdienste) {
			JsonText wetterdaten = wetterdienst.sendeDaten();
			System.out.println(wetterdaten.toString());
		}
	}
} 		    
 		\end{lstlisting}
 	\end{exampleblock}
\end{frame}

\frameWithImageOnly{\BILDERPFAD/adapter1.jpg}
\note{
    Das nächste Beispiel soll zwei inkompatible Systeme miteinander verbinden. Wir haben einen Temperatursensor, der die Temperatur in Fahrenheit misst, wir benötigen jedoch Celsius. Der Fahrenheitsensor kann nicht einfach ersetzt werden, sondern wir müssen eine Möglichkeit finden ihn zu adaptieren und die Werte umzurechnen in Celsius.
}

\begin{frame}{Klassendiagramm - Beispiel}
    \begin{tikzpicture}[scale=0.7, transform shape]
        \umlclass[x=-3,y=0]{Client}{
        }{
        }
    
        \umlclass[x=8,y=-3]{TemperatureAdapter}{
        }{
        
        }

        \umlclass[type=interface, x=2,y=0]{SmartHomeTemperatureSensor}{
        }{
        }

        \umlclass[x=8,y=3]{TemperatureSensor}{
        }{
        }
        
        \umluniassoc{Client}{SmartHomeTemperatureSensor}
        
        \umlinherit{TemperatureSensor}{SmartHomeTemperatureSensor}
        \umlinherit{TemperatureAdapter}{SmartHomeTemperatureSensor}
    \end{tikzpicture}
\end{frame}

\begin{frame}[fragile]{Ziel}
	\begin{exampleblock}{SmartHomeTemperatureSensor.java}
 		\begin{lstlisting}
public interface SmartHomeTemperatureSensor {
	double getTemperature();
}

 		\end{lstlisting}
 	\end{exampleblock}
\end{frame}

\begin{frame}[fragile]{AdaptierteKlasse}
	\begin{exampleblock}{TemperatureSensor}
 		\begin{lstlisting}
public class TemperatureSensor {
	private static final double MAX_FAHRENHEIT = 1000.0;
	
	private Random random;
	
	public TemperatureSensor() {
		random = new Random();
	}
	
	public double getTemperatureInFahrenheit() {
		return random.nextDouble() * MAX_FAHRENHEIT;
	}
}
	    
 		\end{lstlisting}
 	\end{exampleblock}
\end{frame}

\begin{frame}[fragile]{Adapter}
	\begin{exampleblock}{TemperatureAdapter.java}
 		\begin{lstlisting}
public class TemperatureAdapter implements
        SmartHomeTemperatureSensor {
	private TemperatureSensor temperatureSensor;

	public TemperatureAdapter(TemperatureSensor temperatureSensor) {
		this.temperatureSensor = temperatureSensor;
	}

	@Override
	public double getTemperature() {
		double fahrenheit = temperatureSensor.
            getTemperatureInFahrenheit();
		return (fahrenheit - 32) * 5.0 / 9.0;
	}
}    
 		\end{lstlisting}
 	\end{exampleblock}
\end{frame}